\originalchapter{8.05 2013期末：谱生成与纠缠}
原件为Fall 2013 Final Exam, 2013年12月17日9:00--12:00, 六题100分, 闭卷禁计算器. 以下全小问为\textbf{AI独立解答}. 能级图按题意大幅绘制, 图上微小劈裂作示意放大, 数值偏移由表格精确给出.
\originalsection{第1题：十个判断（20分）}
\begin{exercise}
判断: (1)$\sigma_1\otimes\sigma_1$与$\sigma_3\otimes\sigma_1$对易; (2)$\sigma_1\otimes\sigma_1$与$\sigma_3\otimes\sigma_3$对易; (3)$(T\otimes S)^\dagger=S^\dagger\otimes T^\dagger$; (4)有限维单射自映射可逆; (5)角动量守恒禁止自旋$1/2$衰变为两个自旋$1/2$; (6)$[L^2,x_i]=0$; (7)$\boldsymbol r\cdot\boldsymbol p=\boldsymbol p\cdot\boldsymbol r-3\ii\hbar$; (8)旋转矢量$\boldsymbol A$满足$\boldsymbol A\cdot\boldsymbol L=\boldsymbol L\cdot\boldsymbol A$; (9)氢Runge--Lenz矢$\boldsymbol R$本身满足角动量代数; (10)经典和量子都有$\boldsymbol R\cdot\boldsymbol L=0$.
\end{exercise}
\begin{solution}
依次\textbf{F,T,F,T,T,F,F,T,F,T}. (1)只有第一因子反对易, 整体不对易; (2)两个因子各反对易, 双负抵消; (3)伴随保持因子顺序,$T^\dagger\otimes S^\dagger$; (4)秩为维数; (5)末态合成整数不能匹配半整数; (6)位置是矢量算符, 可改变$\ell$, 例如$x_iY_{00}$产生$\ell=1$角依赖; (7)正确为$\boldsymbol r\cdot\boldsymbol p=\boldsymbol p\cdot\boldsymbol r+3\ii\hbar$; (8)$\sum_i[L_i,A_i]=\ii\hbar\sum_{ik}\epsilon_{iik}A_k=0$; (9)$[R_i,R_j]$与$HL_k$成比例, 不是$\ii\hbar\epsilon_{ijk}R_k$; (10)见上一期末第3题量子算符证明.
\end{solution}
\originalsection{第2题：位移压缩态的占有数（10分）}
\begin{exercise}
$\ket{\alpha,\gamma}=D(\alpha)S(\gamma)\ket0$,$\alpha$复、$\gamma$实, 求$\langle N\rangle$并判断自由谐振子演化中是否改变.
\end{exercise}
\begin{solution}
$D^\dagger aD=a+\alpha$, 压缩真空$\langle a\rangle=0$, $\langle a^\dagger a\rangle=\sinh^2\gamma$. 逐项展开$D^\dagger ND=(a^\dagger+\alpha^*)(a+\alpha)$, 得$\langle N\rangle=|\alpha|^2+\sinh^2\gamma$. $[H,N]=0$所以恒定. 分别代表位移的经典振动能量与压缩增加的真空涨落能量, 不随椭圆噪声方向旋转而变.
\end{solution}
\originalsection{第3题：球方阱束缚门槛（10分）}
\begin{exercise}
$V=-V_0$（$r<a$）、0（$r>a$）,$V_0>0$. 推导至少有一个三维束缚态所需深度不等式.
\end{exercise}
\begin{solution}
最容易束缚的是$\ell=0$, 因所有$\ell>0$多了非负离心势. 正规径向$u(0)=0$, 内$A\sin kr$, 外$B\e^{-\kappa r}$, 匹配$k\cot ka=-\kappa$, $k^2+\kappa^2=2mV_0/\hbar^2$. 第一束缚根要$z_0=a\sqrt{2mV_0}/\hbar>\pi/2$, 即$V_0>\pi^2\hbar^2/(8ma^2)$. 等号零能外部$u$为常数, $\psi=u/r$虽点值衰减却$\int|u|^2\dd r$发散, 不算正规束缚态. 与一维全线任意浅吸引方阱有束缚态不同, 三维有深度门槛.
\end{solution}
\originalsection{第4题：一个谱生成角动量（15分）}
\begin{exercise}
本题$\hbar=1$, $j$正整数,$J_z=-j+N$,$J_+=a^\dagger\sqrt{2j-N}$,$J_-=\sqrt{2j-N}a$. (a)在平方根合法态上验证代数. (b)求$J^2$及对易. (c)构造$j$多重态并检验作用. (d)与振子$H=N+1/2$是否对易.
\end{exercise}
\begin{solution}
在$0\leq n\leq2j$有限不变空间上,$J_+\ket n=\sqrt{(n+1)(2j-n)}\ket{n+1}$,$J_-\ket n=\sqrt{n(2j-n+1)}\ket{n-1}$. 端点向外系数0, 所以不需对负平方根运算. $J_z$作用$n-j$, 立即$[J_z,J_\pm]=\pm J_\pm$. $J_+J_-$系数$n(2j-n+1)$,$J_-J_+$系数$(n+1)(2j-n)$, 差$2(n-j)$, 因而$[J_+,J_-]=2J_z$. 用$J^2=J_z^2+(J_+J_-+J_-J_+)/2$得$j(j+1)I$, 与三分量对易. 态$\ket{j,m}=\ket{n=j+m}$,$m=-j,\ldots,j$, 给预期阶梯系数$\sqrt{j(j+1)-m(m\pm1)}$. $[H,J_z]=0$, $[H,J_\pm]=\pm J_\pm$, 后二者改变振子能量, 因而是\textbf{谱生成}而非同能量旋转对称. 对全无限振子空间没有实平方根的该定义, 必须保留有限子空间限制.
\end{solution}
\originalsection{第5题：四维振子的简并劈裂（20分）}
\begin{exercise}
四维等频$H=\hbar\omega(N+2)$. (a)列$N=0,1,2$态并计数. (b)把前两个坐标、后两个坐标各组成二维振子, 数$n_a,n_b$对应Schwinger角动量$j_1=n_a/2,j_2=n_b/2$, 列$N=0,1,2,3$的$j_1\otimes j_2$及维数. (c)加$2E_0\boldsymbol J^{(1)}\cdot\boldsymbol J^{(2)}/\hbar^2$,$0<E_0\ll\hbar\omega$, 求四组劈裂、表示与态数, 画大幅能级图.
\end{exercise}
\begin{solution}
$N=0$真空1态; $N=1$为四个$a_i^\dagger\ket0$; $N=2$为四个$(a_i^\dagger)^2\ket0$及六个$a_i^\dagger a_j^\dagger\ket0$（$i<j$）, 共10. 一般简并$\binom{N+3}3$. 各二维振子的两占有数在固定总数$n$形成$j=n/2$, 维数$n+1$. 表如下:
\begin{center}\begin{tabular}{ccc}\toprule $N$&$j_1\otimes j_2$&各表示维数\\\midrule
0&$0\otimes0$&1\\
1&$\tfrac12\otimes0,\ 0\otimes\tfrac12$&2,2\\
2&$1\otimes0,\ \tfrac12\otimes\tfrac12,\ 0\otimes1$&3,4,3\\
3&$\tfrac32\otimes0,\ 1\otimes\tfrac12,\ \tfrac12\otimes1,\ 0\otimes\tfrac32$&4,6,6,4\\\bottomrule\end{tabular}\end{center}
耦合$\boldsymbol J=\boldsymbol J^{(1)}+\boldsymbol J^{(2)}$, 偏移精确为$\delta E=E_0[j(j+1)-j_1(j_1+1)-j_2(j_2+1)]$. 因为扰动与原$H,J_1^2,J_2^2,J^2$对易, 这个线性偏移在所给模型为精确结果, 不仅是第一阶.

$N=0$一态无偏移. $N=1$两套$j=1/2$, 共4态无偏移. $N=2$: 两套$j=1$来自$1\otimes0,0\otimes1$, 共6态无偏移; $\tfrac12\otimes\tfrac12$分成$j=1$三态偏$+E_0/2$, $j=0$一态偏$-3E_0/2$. $N=3$: 两套$j=3/2$来自端点表示, 共8态无偏移; 两个中间表示各含$j=3/2$四态偏$+E_0$与$j=1/2$两态偏$-2E_0$, 总8、4态. 
\clearpage
以下图中括号为态数, 劈裂相对主间隔放大, 所有表示来源已在文字和表中逐一列出.
\begin{center}\begin{tikzpicture}[x=1cm,y=1.3cm,>=Stealth]
\draw[->](0,-.5)--(0,10.6)node[above]{$E$};
\foreach \y/\lab/\num in {0/{2\hbar\omega}/1,3/{3\hbar\omega}/4,6/{4\hbar\omega}/10,9/{5\hbar\omega}/20}{
\draw(.5,\y)--(1.6,\y);\node[left]at(-.15,\y){$\lab$};\node[above]at(1,\y){$N=\pgfmathparse{int(\y/3)}\pgfmathresult\ (\num)$};}
\draw[main](3,0)--(5,0)node[right]{$j=0$ (1)};\draw[->](1.7,0)--(2.9,0);
\draw[main](3,3)--(5,3)node[right]{$j=1/2$ (4)};\draw[->](1.7,3)--(2.9,3);
\draw[main](3,6)--(5,6)node[right]{$\delta E=0,\ j=1$ (6)};
\draw[main](3,6.35)--(5,6.35)node[right]{$+E_0/2,\ j=1$ (3)};
\draw[main](3,4.95)--(5,4.95)node[right]{$-3E_0/2,\ j=0$ (1)};
\foreach \y in {4.95,6,6.35}{\draw[->](1.7,6)--(2.9,\y);}
\draw[main](3,9)--(5,9)node[right]{$\delta E=0,\ j=3/2$ (8)};
\draw[main](3,9.7)--(5,9.7)node[right]{$+E_0,\ j=3/2$ (8)};
\draw[main](3,7.6)--(5,7.6)node[right]{$-2E_0,\ j=1/2$ (4)};
\foreach \y in {7.6,9,9.7}{\draw[->](1.7,9)--(2.9,\y);}
\node at(4,-.6){原组$\longrightarrow$耦合后（小劈裂放大示意）};
\end{tikzpicture}\end{center}
\end{solution}
\clearpage
\originalsection{第6题：纠缠的产生与局部基变换（25分）}
\begin{exercise}
$H=\hbar\omega\sigma_3\otimes\sigma_3$, $Z$基$\ket{++},\ket{+-},\ket{-+},\ket{--}$,$\ket{x;\pm}=(\ket+\pm\ket-)/\sqrt2$. (a)写$H,U$矩阵. (b)演化初态$\ket{x;+}\ket{x;+}$, 找无纠缠时刻. (c)把$\exp[\ii(K\otimes I+I\otimes K)]$分为张量积. (d)将任意$\alpha\ket{x;+}\ket{x;+}+\beta\ket{x;-}\ket{x;-}$变为$\alpha\ket{++}+\beta\ket{--}$, 用$H'=H_0\otimes I+I\otimes H_0$和$K=-H_0t'/\hbar$构造.
\end{exercise}
\begin{solution}
$H=\hbar\omega\operatorname{diag}(1,-1,-1,1)$,$U=\operatorname{diag}(\e^{-\ii\omega t},\e^{\ii\omega t},\e^{\ii\omega t},\e^{-\ii\omega t})$. 初态演化为这些四个相位各乘$1/2$. 二比特纯态系数矩阵行列式为$(\e^{-2\ii\omega t}-\e^{2\ii\omega t})/4=-(\ii/2)\sin2\omega t$, 因此无纠缠当$\omega t=k\pi/2$,$k\in\mathbb Z$; $\omega=0$则一直无纠缠. 行列式非零意味着不能写两列向量的外积, 即非乘积态.

$K\otimes I$与$I\otimes K$对易, 故指数为$\e^{\ii K}\otimes\e^{\ii K}$. 令Hadamard矩阵$F=(\sigma_x+\sigma_z)/\sqrt2$, 满足$F\ket{x;+}=\ket+$,$F\ket{x;-}=\ket-$且$F^2=I$. 取Hermitian$K=\pi(I-F)/2$, 在$F$的$\pm1$谱上$K$分别0、$\pi$, 所以$\e^{\ii K}=F$精确成立. 任取$t'>0$,$H_0=-\hbar K/t'$即完成指定转变, 包括相对相位而非仅到模方相等. 局部幺正只改变表示, 不改变两态系数$\alpha,\beta$的Schmidt值或纠缠量; 本题不是用局部操作消除纠缠.
\end{solution}
